\clearpage
\appendix
\section*{Phụ lục}
\addcontentsline{toc}{section}{Phụ lục}
\lstset{
  basicstyle=\ttfamily\scriptsize,
  breaklines=true,
  breakatwhitespace=false,
  columns=fullflexible,
  keepspaces=true,
  showstringspaces=false,
  frame=single,
  xleftmargin=2mm,
  xrightmargin=2mm
}

\section{Prompt zero-shot}

Các đoạn dưới đây là \emph{system prompt} P0--P3 nguyên văn từ prompt registry
của Giai đoạn 2. Câu hỏi và năm context được gửi riêng trong user prompt;
phần này không đưa đáp án mẫu cho generator.

\noindent\textbf{P0 --- Đối chứng.} Yêu cầu trả lời ngắn, dựa vào context và có
citation cho từng phát biểu thực tế.
\begin{quote}\small\sloppy
You are a helpful assistant answering questions based on provided context.
Answer concisely and only based on the given context. Cite supporting context
using its bracketed number, for example [1]. Every factual claim must have a
citation, and you must not cite context that does not support the claim. If the
answer is not in the context, say: 'I cannot find this information in the
provided context.'
\end{quote}

\noindent\textbf{P1 --- Siết grounding.} Kiểm tra liệu cấm suy đoán và chỉ dẫn
bằng chứng trực tiếp có tăng Faithfulness hay không.
\begin{quote}\small\sloppy
Answer the question using only facts explicitly stated in the numbered
context. Do not add background knowledge, assumptions, or plausible details.
Give a concise answer and cite the supporting context with [n] after every
factual claim. Cite only context that directly supports that claim. If the
context does not contain the answer, say: 'I cannot find this information in
the provided context.'
\end{quote}

\noindent\textbf{P2 --- Đúng loại đáp án.} Kiểm tra liệu đáp án một câu, chỉ
chứa loại thông tin được hỏi, có tăng độ chính xác hay không.
\begin{quote}\small\sloppy
Answer using only the numbered context. Give one short answer sentence
containing only the information type requested by the question, such as a
person, place, date, year, number, object, event, or explanation. Include
multiple items only when the question explicitly asks for them. Do not repeat
the question, list alternative answers, or add background details. Place a
supporting citation [n] immediately after the answer. If the answer is not in
the context, say: 'I cannot find this information in the provided context.'
\end{quote}

\noindent\textbf{P3 --- Phân biệt sự kiện.} Kiểm tra liệu chọn đoạn khớp chủ
thể, sự kiện, địa điểm và thời gian có giảm nhầm lẫn giữa các bài hay không.
\begin{quote}\small\sloppy
Answer using only facts explicitly stated in the numbered context. First
identify the passage that best matches all identifying details in the
question, including its subject, event, location, and time. Do not combine
facts, figures, or dates from different events, articles, or reporting
periods. Give one short answer sentence containing only the requested
information type, and include multiple items only when explicitly asked.
Place a supporting citation [n] immediately after the answer and cite only
the passage that directly supports it. If no passage contains a defensible
answer, say: 'I cannot find this information in the provided context.'
\end{quote}

\section{Prompt one-shot P2-1S}

P2-1S giữ nguyên P2 và nối thêm ví dụ dưới đây vào cùng system prompt. Ví dụ
\textbf{hoàn toàn được tạo tổng hợp}: bảo tàng Northbridge, kế hoạch Riverside,
ngữ cảnh, câu hỏi và đáp án đều không lấy từ NewsQA, corpus truy xuất,
development hoặc held-out. Ví dụ chỉ minh họa cách trả lời ngắn và trỏ tới đúng
số thứ tự context.

\begin{quote}\small
Demonstration (format example only):\\
Context [1]: The Northbridge museum opened a new gallery in 2018.\\
Context [2]: The city council approved the Riverside transit plan in March
2021.\\
Question: When was the Riverside transit plan approved?\\
Answer: March 2021 [2]

Now answer the user's question using the supplied numbered context and the
same concise citation format.
\end{quote}

Bản prompt đầy đủ là P2 ở Phụ lục A, hai dòng trống, rồi phần demonstration
trên. Bản ghi nguyên văn và SHA-256 được lưu trong
\path{results/phase2/1s_prompt_experiment/1s_dev_results/results/one_shot_prompt_contract.json}.

\section{Prompt judge RAGAS}

Các file dưới đây được xuất bằng \path{scripts/export_report_ragas_prompts.py}
từ \textbf{RAGAS 0.4.3}. Mỗi file chứa instruction nguyên văn, JSON output
schema, toàn bộ ví dụ few-shot có sẵn trong thư viện và phần input cuối. Các
giá trị trong dấu \texttt{<...>} ở input cuối là \emph{giữ chỗ} để in template;
khi chấm, RAGAS thay bằng dữ liệu của từng câu hỏi. Judge là GLM-5.3-Flash qua
Fireworks, reasoning \texttt{low}. Citation metrics không sử dụng judge này.

\begin{table}[H]
  \centering
  \caption{Template RAGAS được dùng cho năm metric.}
  \small
  \begin{tabularx}{\textwidth}{@{}p{38mm}X@{}}
    \toprule
    \textbf{Metric} & \textbf{Template / bước chấm} \\
    \midrule
    Faithfulness & Tách đáp án thành phát biểu; kiểm tra từng phát biểu với context. \\
    Answer Correctness & Tách đáp án và gold thành phát biểu; phân loại TP/FP/FN. \\
    Answer Relevancy & Sinh câu hỏi từ đáp án, phát hiện trả lời lảng tránh; so embedding với câu hỏi gốc. \\
    Context Precision & Chấm từng context có hữu ích cho câu hỏi và gold không; xét thứ tự context. \\
    Context Recall & Chấm từng phát biểu của gold có được context hỗ trợ không. \\
    \bottomrule
  \end{tabularx}
\end{table}

\subsection*{Tách phát biểu: Faithfulness và Answer Correctness}
\lstinputlisting{appendix/ragas_prompts/statement_generation.txt}

\subsection*{Kiểm tra phát biểu: Faithfulness}
\lstinputlisting{appendix/ragas_prompts/faithfulness_verdict.txt}

\subsection*{Phân loại TP/FP/FN: Answer Correctness}
\lstinputlisting{appendix/ragas_prompts/answer_correctness_classification.txt}

\subsection*{Sinh câu hỏi: Answer Relevancy}
\lstinputlisting{appendix/ragas_prompts/answer_relevancy.txt}

\subsection*{Chấm từng context: Context Precision}
\lstinputlisting{appendix/ragas_prompts/context_precision.txt}

\subsection*{Kiểm tra độ phủ gold: Context Recall}
\lstinputlisting{appendix/ragas_prompts/context_recall.txt}

\section{Cách chấm citation}

Generator nhận các context đánh số \texttt{[1]}--\texttt{[5]}. Bộ chấm ánh xạ
chỉ số trong đáp án về chunk ID theo đúng thứ tự context đã gửi. Ví dụ tổng
hợp: \texttt{[1]}--\texttt{[5]} lần lượt là \texttt{c1}--\texttt{c5}, còn tập
gold là \texttt{\{c2,c4\}}.

\begin{table}[H]
  \centering
  \caption{Ví dụ phân biệt citation hợp lệ và citation đúng bằng chứng.}
  \label{tab:appendix-citation-examples}
  \small
  \begin{tabularx}{\textwidth}{@{}p{24mm}p{28mm}rrrX@{}}
    \toprule
    \textbf{Citation} & \textbf{Chunk được cite} & \textbf{Validity} &
    \textbf{Precision} & \textbf{Recall} & \textbf{Giải thích} \\
    \midrule
    \texttt{[2]} & \texttt{c2} & 1 & 1 & 0,5 & Hợp lệ, đúng một trong hai gold. \\
    \texttt{[3]} & \texttt{c3} & 1 & 0 & 0 & Hợp lệ nhưng sai bằng chứng. \\
    \texttt{[6]} & Không có & 0 & 0 & 0 & Vượt số context đã gửi. \\
    \texttt{[2][3]} & \texttt{c2,c3} & 1 & 0,5 & 0,5 & Có một citation thừa. \\
    \bottomrule
  \end{tabularx}
\end{table}

\begin{itemize}
  \item \textbf{Citation Validity} $= |C|/(|C|+|I|)$, với $C$ là tập chunk ID
  được cite từ chỉ số hợp lệ và $I$ là tập chỉ số citation ngoài phạm vi.
  Không có citation thì điểm bằng 0.
  \item \textbf{Citation Precision} $=|C\cap G|/|C|$: bao nhiêu chunk được
  cite thực sự thuộc gold $G$.
  \item \textbf{Citation Recall} $=|C\cap G|/|G|$: tỷ lệ gold chunks được cite.
  \item \textbf{Citation F1} $=2PR/(P+R)$ khi $P+R>0$; nếu không thì bằng 0.
  \item Bộ chấm loại citation trùng trước khi tính và tính điểm từng câu rồi
  lấy trung bình. Validity chỉ kiểm tra chỉ số tồn tại, không xác nhận claim
  được context hỗ trợ về mặt ngữ nghĩa.
\end{itemize}

\section{Nguồn tái lập}

\begin{itemize}
  \item Prompt generator: \path{configs/experiments/phase2_generation_prompts.yaml};
  prompt one-shot và SHA-256: \path{one_shot_prompt_contract.json} trong artifact
  development P2-1S.
  \item Prompt judge: các file \texttt{.txt} và \texttt{manifest.json} trong
  \path{docs/latex/report/appendix/ragas_prompts/}; chạy lại
  \path{scripts/export_report_ragas_prompts.py} với RAGAS 0.4.3 để đối chiếu.
  \item Bộ chấm citation: \path{common/newsqa_rag/evaluation/metrics.py}; ánh xạ
  \texttt{[n]} và loại chỉ số ngoài phạm vi: \path{common/newsqa_rag/response_parsing.py}.
  \item Mỗi \texttt{run\_manifest.json} ghi model, question IDs, prompt hash và
  coverage. Bảng điểm đầy đủ ở \path{results/phase2/}.
\end{itemize}
